\section{可选的有界多起点轨迹重启与优化}
\label{sec:geometry-refinement}

旧几何修正先进行启发式 PDDP 更新，再用终端投影改善终点。这可以产生可行解，但不能由此推断目标函数达到驻点。主迭代采用实际 merit 回溯和拒绝回滚；终端投影仍只降低终端误差，不优化全程控制变化。正式避障示例现在显式启用一个独立的有界多起点射击优化阶段；关闭射击阶段的入口保留；该组合不是标准 CBF 或文献算法的逐式复现。

设 $N$ 为状态节点数，$B_{kj}$ 为 $K$ 个等距控制节点形成的分段线性基函数，优化变量是控制节点 $z_j$ 与可选时间 $T$：
\[
 u_k=\sum_{j=1}^{K}B_{kj}z_j,\qquad
 x_{k+1}=x_k+\frac{T}{N-1}f(x_k,u_k).
\]
先优化控制再重新积分状态，最终控制仍按零阶保持应用；没有对输出位置曲线插值或平滑。节点控制满足幅值界，凸组合后的采样控制也满足同一幅值界。当前阶段仅支持具有解析动力学导数的二次代价模型。

原二次目标 $J$ 保持不变并单独报告。重启阶段优化的辅助目标包含控制变化项：
\[
 J_{\mathrm{reg}}=J+\frac{\lambda}{2}
  \sum_{k=1}^{N-2}\frac{\|u_{k+1}-u_k\|^2}{h},\qquad h=\frac{T}{N-1}.
\]
为改善终点和几何可行性，非线性最小二乘 merit 另外包含
\[
 \frac{10^5}{2}\|x_N-x_f\|^2+
 \frac{10^5}{2}\sum_{o,k}\max(0,\delta-c_{o,k})^2,
 \qquad \delta=0.005\ \mathrm{m},
\]
其中 $c_{o,k}$ 是状态线段到物理障碍半径加安全余量的精确净空。有限罚权不是严格约束证明；所有候选在最终选择时还必须满足原有终端误差、控制界及精确线段净空阈值。

状态对控制节点和时间的灵敏度通过动力学链式传播。Gauss--Newton 方向由
\[
 (A^\top A+\mu D)d=-A^\top r,\qquad
 D=\operatorname{diag}(\max(1,\operatorname{diag}(A^\top A)))
\]
求得，采用 Cholesky 分解，$\mu$ 最小为 $10^{-5}$。控制和时间候选裁剪到其边界，只有实际 merit 严格降低才接受，否则增加阻尼后重算方向。罚函数下降历史与原 PDDP 试探历史分别保存。\code{geometric_refinement_feasible} 表示已选出改善且可行的解，不表示已证明全局最优或满足严格 KKT 驻点条件。

二维使用正、负两个余弦转向控制初值，三维分别沿两路控制使用正、负初值，共四个。它们独立于历史 PDDP 轨迹的绕行方向，避免放大旧迭代的舍入敏感性；近似并列的 merit 使用固定次序选择。
候选与基线均通过公共数值验收检查物理可行性，包括状态约束、控制界、时间界、Euler/RK4 缺陷以及适用的解析全程几何检查。
若要求代价下降，候选原目标不得超过原始初值目标（相对容差 $10^{-10}$）；该参照不是可能不可行的廉价投影基线。
通过可行性和代价策略的候选可替换不合格基线，即使其目标更大；合格可行解之间按 $J_{\mathrm{reg}}$ 的既有改善阈值选择。
没有合格候选则保留原输出并继续正常验收，不能因此标记成功。
各起点的次数、目标、控制变化项、终端误差和净空均保存在原始记录中；正式轨迹图展示最终选择的解。

正式示例参数为 \code{geometric_refinement_iterations=80}、\code{geometric_refinement_knots=12}、\code{geometric_control_rate_weight=0.02} 和 \code{geometric_refinement_free_time=true}；接口默认迭代数为 0，便于保留原求解行为。设置对应字段为 0 可禁用新阶段。新示例时间界为 $5\le T\le6$ s。

\subsection{解析保持控制回放与区间净空下界}
二维匀速转向和三维恒速航迹模型的每段零阶保持运动可解析积分。
例如
\[
 \int_0^h\cos(a+bt)\,dt=h\,\operatorname{sinc}(bh/2)\cos(a+bh/2),
\]
此处 $\operatorname{sinc}(z)=\sin(z)/z$；小角采用级数避免除零。
三维水平速度用积化和差公式积分，航迹角和航向角按各自角速度推进。
回放从原初态连续推进，不在每个 Euler 节点重新起步。

每段控制分成 64 个子区间，计算解析位置到所有静态圆/球障碍的净空。
距离函数关于位置为 1-Lipschitz，速度幅值为 $v_{\max}$，因此区间内净空下界为
\[
 c_{\rm lower}=\min_i c(x(t_i))-v_{\max}\Delta t_{\max}/2
 -10^{-10}(1+\|x_0\|+v_{\max}T).
\]
最后一项为浮点计算预留裕量；这不是严格区间算术证明，也不涵盖模型误差或扰动。
MATLAB 的 \code{ddpack.verifyGeometricRollout} 与 C++ 的独立诊断均计算该下界，
适用模型必须下界非负且解析终端误差满足原容差，才能标为成功。
其他动力学的该指标为不适用，不能将空值解释为通过此检查。
Python 第三套实现同时核对解析回放与 RK4 的逐点一致性。
对本单圆/单球阻挡场景，可用忽略起终航向及控制限制的几何最短绕障距离作参照。沿起终连线的球心到两端距离为 $d_1,d_2$，膨胀半径为 $R$，其长度为
\[
 L_{\mathrm{geom}}=\sqrt{d_1^2-R^2}+\sqrt{d_2^2-R^2}
 +R\left(\pi-\arccos\frac{R}{d_1}-\arccos\frac{R}{d_2}\right).
\]
二维与三维分别约为 5.02502 m 和 5.49492 m；新返回路径约为 5.03221 m 和 5.50032 m，只分别多约 0.14\% 和 0.10\%。终点仍按给定误差容差验收。该长度参照不证明原二次目标的全局最优性，也不排除其他参数化获得更低目标。
